% ---------------------------------------------------------------------
\chapter{Warum Versionsverwaltung?}
\label{kap:warum}

\begin{lernziele}
\begin{itemize}[nosep]
  \item die Probleme kennen, die ohne Versionsverwaltung entstehen,
  \item verstehen, was ein Versionsverwaltungssystem leistet,
  \item Git und GitHub klar voneinander unterscheiden,
  \item wissen, warum Git in der Industrie zum Standard geworden ist.
\end{itemize}
\end{lernziele}

\section{Das Problem}

Jeder kennt Ordner, die so aussehen:

\begin{lstlisting}[style=ausgabe]
bericht.docx
bericht_v2.docx
bericht_v2_korrigiert.docx
bericht_final.docx
bericht_final_wirklich_final.docx
bericht_final_wirklich_final_Anna.docx
\end{lstlisting}

Bei einem einzelnen Dokument ist das lästig. Bei Software mit Hunderten von Dateien und
mehreren Beteiligten wird es unbeherrschbar. Vier Fragen lassen sich ohne Werkzeug kaum
noch beantworten:

\begin{enumerate}
  \item \textbf{Was hat sich geändert?} Welche Zeilen sind zwischen gestern und heute anders?
  \item \textbf{Wer hat es geändert und warum?} Wer hat die Zeile eingefügt, die seit Montag
        einen Fehler verursacht, und was war die Absicht?
  \item \textbf{Wie komme ich zurück?} Gestern lief das Programm noch. Wie stelle ich den
        damaligen Zustand wieder her, ohne die Arbeit von heute zu verlieren?
  \item \textbf{Wie arbeiten mehrere Personen gleichzeitig?} Anna und Ben ändern am selben
        Tag dieselbe Datei. Wessen Version gilt?
\end{enumerate}

\section{Die Lösung: ein Versionsverwaltungssystem}

\begin{definition}{Versionsverwaltung}
Ein \begriff{Versionsverwaltungssystem} (englisch \emph{Version Control System}, VCS)
speichert die Geschichte eines Projekts als Folge von Zwischenständen. Zu jedem
Zwischenstand werden Autor, Zeitpunkt und eine Beschreibung festgehalten. Jeder frühere
Stand lässt sich jederzeit wiederherstellen und mit anderen vergleichen.
\end{definition}

Damit sind die vier Fragen beantwortet: Das System weiss, \emph{was} sich geändert hat,
\emph{wer} es war und \emph{warum}, es kann zu jedem Stand \emph{zurück}, und es kann die
Arbeit mehrerer Personen \emph{zusammenführen}.

Ein Versionsverwaltungssystem ersetzt dabei kein Backup und keinen Cloud-Speicher wie
OneDrive. Der Unterschied liegt in der Absicht: Ein Cloud-Speicher synchronisiert
\emph{automatisch} jede Änderung. Bei der Versionsverwaltung entscheidet der Mensch
\emph{bewusst}, wann ein Stand gut genug ist, um festgehalten zu werden, und beschreibt ihn.
Dadurch entsteht eine lesbare Geschichte statt eines Stroms von Dateiversionen.

\section{Git}

\begriff{Git} ist das mit Abstand verbreitetste Versionsverwaltungssystem. Es ist kostenlos,
quelloffen und läuft auf jedem Betriebssystem.

\begin{hintergrund}
Linus Torvalds, der Erfinder von Linux, entwickelte Git im Jahr 2005 in wenigen Wochen.
Das Linux-Projekt hatte damals Tausende Mitwirkende auf der ganzen Welt und verlor über
Nacht das Recht, sein bisheriges Werkzeug kostenlos zu nutzen. Torvalds' Anforderungen waren
entsprechend hoch: Git sollte sehr schnell sein, auch ohne Internetverbindung funktionieren,
und kein einzelner Server sollte das Projekt lahmlegen können. Diese Entscheidungen prägen
Git bis heute.
\end{hintergrund}

Die wichtigste dieser Entscheidungen: Git ist \begriff{verteilt}. Ältere Systeme hatten
einen zentralen Server, auf dem die Geschichte lag. Ohne Verbindung zum Server konnte man
nicht arbeiten. Bei Git besitzt dagegen \textbf{jede Kopie die vollständige Geschichte}.
Man kann offline Zwischenstände speichern, die Geschichte durchsuchen und Versionen
vergleichen. Synchronisiert wird nur, wenn man es will.

\section{Git und GitHub: zwei verschiedene Dinge}

Einsteiger verwechseln die beiden häufig, dabei ist die Trennung einfach:

\begin{center}
\renewcommand{\arraystretch}{1.35}
\begin{tabularx}{\linewidth}{@{}l>{\raggedright\arraybackslash}X>{\raggedright\arraybackslash}X@{}}
  \toprule
                     & \textbf{Git} & \textbf{GitHub} \\
  \midrule
  Was es ist         & ein Programm auf dem eigenen Computer & ein Internetdienst (Website) \\
  Aufgabe            & Geschichte speichern, vergleichen, zusammenführen & Kopien im Internet aufbewahren und teilen \\
  Zusätzliches       & –  & Pull Requests, Code Review, automatische Tests, Fehlerverwaltung \\
  Gehört zu          & freie Software, keiner Firma & Microsoft (seit 2018) \\
  Alternativen       & –  & GitLab, Bitbucket, eigener Server \\
  \bottomrule
\end{tabularx}
\end{center}

Ein Vergleich hilft: Git ist wie die Fotokamera, GitHub wie das Online-Fotoalbum, in dem
man Bilder sichert und mit anderen teilt. Die Kamera funktioniert auch ohne Album. Das Album
ist aber der Ort, an dem die Zusammenarbeit stattfindet.

\begin{praxis}
In Unternehmen ist die Plattform (GitHub, GitLab oder Bitbucket) das Zentrum der
Entwicklungsarbeit. Dort wird nicht nur Code aufbewahrt: Jede Änderung wird dort von
Kolleginnen und Kollegen geprüft (\emph{Code Review}), automatisch getestet und erst dann
übernommen. Diese Abläufe werden in Kapitel~\ref{kap:pr} behandelt.
\end{praxis}

\begin{zusammenfassung}
\begin{itemize}[nosep]
  \item Versionsverwaltung beantwortet: Was hat sich geändert, wer war es und warum,
        wie komme ich zurück, wie arbeiten mehrere Personen zusammen?
  \item Git ist das Werkzeug auf dem eigenen Computer, GitHub der Dienst im Internet.
  \item Git ist verteilt: Jede Kopie enthält die ganze Geschichte.
\end{itemize}
\end{zusammenfassung}
